% Generated human-review packet template.  Source records, Lean statements,
% and raw-source-to-Spec screening fields are substituted by
% scripts/review_dashboard_packet.py.
% Do not hand-edit generated packets; annotate the PDF or reconcile notes into
% the paper-local human-review log.
\documentclass[10pt]{article}
\usepackage[margin=0.43in,footskip=0.2in]{geometry}
\usepackage{fontspec}
% Use a Unicode-complete main face as well as a Unicode-complete code face.
% Reviewer reasons and source labels can contain exact Greek symbols outside
% verbatim blocks; silently dropping one would corrupt the review packet.
\setmainfont{DejaVu Serif}
% The broad DejaVu Sans Unicode coverage preserves Lean's mathematical glyphs
% (including U+2216 set difference and superscript identifiers) in verbatim
% review blocks.  Its proportional metrics are preferable to silently missing
% exact semantic text from a narrower monospace font.
\setmonofont{DejaVu Sans}
\usepackage{hyperref}
\usepackage{amssymb}
\usepackage{fvextra}
\usepackage{enumitem}
\setlength{\parindent}{0pt}
\setlength{\parskip}{0.15em}
\linespread{0.92}
\hypersetup{colorlinks=true,linkcolor=blue,urlcolor=blue,breaklinks=true}
\DefineVerbatimEnvironment{ReviewVerbatim}{Verbatim}{breaklines=true,breakanywhere=true,fontsize=\scriptsize}
\newcommand{\reviewerbox}{%
  \par\noindent\fbox{\parbox{0.96\linewidth}{\rule{0pt}{0.38in}}}\par
}
\newcommand{\reviewmatch}[1]{%
  \par\noindent\textbf{Reviewer check:} $\square$\; #1\par
  \noindent\emph{If left unchecked, explain the discrepancy or uncertainty below.}\par
}

\begin{document}
\fontsize{9}{10}\selectfont

\begin{center}
{\LARGE Human review packet: MBJG25ProducerFairness}\\[0.35em]
\small Generated from byte-pinned source excerpts, the expanded PaperInterface
specifications, and hash-bound raw-source-to-Spec screening records.
\end{center}

\noindent\textbf{Paper:} MBJG25ProducerFairness\\
\textbf{Paper id:} \texttt{MBJG25ProducerFairness}\\
\textbf{Source version:} \parbox[t]{0.76\linewidth}{\raggedright ICWSM 2025, 19(1), 1139-1157}\\
\textbf{Source record:} \texttt{audit/paper\_statement\_map.json}\\
\textbf{Rows:} 2\\
\textbf{Generated:} 2026-09-06\\
\textbf{Source URL:} \href{https://ojs.aaai.org/index.php/ICWSM/article/view/35865}{official source}





\section*{How to use this packet}
For each source claim, compare the exact source excerpts with the expanded
PaperInterface specification. Mark up this PDF or fill the blank reviewer box in the TeX source;
return the annotated file for reconciliation into the paper-local human-review
log.

\section*{Open the interactive dashboard}
The interactive dashboard is an optional alternative to using this PDF; it
presents the same review surface rather than an additional review requirement.
After forking and cloning (or pulling) AppliedModelingLib, run the following from the
repository root. The cache command is needed only the first time in a new clone.
\begin{ReviewVerbatim}
cd <your-repository-checkout>
lake exe cache get
python3 scripts/review_dashboard.py --paper MBJG25ProducerFairness --serve
\end{ReviewVerbatim}
Open \url{http://127.0.0.1:8765/} in a browser and keep that terminal running
while reviewing. The dashboard records saved annotations in this paper's local
reviewer-owned review record.

\section*{Contents}
\begin{itemize}[leftmargin=1.2em]
\item \hyperlink{paper-prerequisite-section-5ef5ef0364b6939c}{Paper-specific semantic prerequisites (1; not paper claims)}
\begin{itemize}[leftmargin=1.2em]
\item \hyperlink{paper-prerequisite-a998c700d4f94308}{MBJG25\allowbreak{}Producer\allowbreak{}Fairness.\allowbreak{}paper\_\allowbreak{}fixed\_\allowbreak{}binary\_\allowbreak{}model\_\allowbreak{}formula\allowbreak{}Spec}
\end{itemize}
\item \hyperlink{source-section-f10b5f68e4acf3dc}{Source claims (2)}
\begin{itemize}[leftmargin=1.2em]
\item \hyperlink{source-claim-98cc132a8c76d666}{paper\_\allowbreak{}theorem3\_\allowbreak{}1\allowbreak{}Spec}
\item \hyperlink{source-claim-ad27f52afb94912c}{paper\_\allowbreak{}theorem3\_\allowbreak{}2\allowbreak{}Spec}
\end{itemize}
\end{itemize}

\clearpage
\hypertarget{paper-prerequisite-section-5ef5ef0364b6939c}{}
\section*{Paper-specific semantic prerequisites}
\hypertarget{paper-prerequisite-a998c700d4f94308}{}
\subsection*{\small\ttfamily\raggedright MBJG25\allowbreak{}Producer\allowbreak{}Fairness.\allowbreak{}paper\_\allowbreak{}fixed\_\allowbreak{}binary\_\allowbreak{}model\_\allowbreak{}formula\allowbreak{}Spec}
\noindent\textbf{Paper-local declaration:}\par
{\footnotesize\ttfamily\raggedright MBJG25\allowbreak{}Producer\allowbreak{}Fairness.\allowbreak{}paper\_\allowbreak{}fixed\_\allowbreak{}binary\_\allowbreak{}model\_\allowbreak{}formula\allowbreak{}Spec\par}
\textbf{Source locator:} {\footnotesize\raggedright cited publication, lines 159-\allowbreak{}219\par}
\paragraph{Verbatim paper-source connection}
\begin{ReviewVerbatim}
with binary reviews, our model extends readily to ordinal                       is the number of negative ratings). We can compute the pos-
reviews as well (see Appendix E).                                               terior mean estimated quality for product v at timestep t as:
                                                                                                                    α̂ + R(v, t)
2.1    Model Primitives                                                                       q̂α̂,β̂ (v, t) =                         .   (1)
Our modeling framework is as follows.                                                                            α̂ + β̂ + |S(v, t)|
   Products. There is a universe of products V . Product v ∈                       The platform impacts the operation of a prior-weighted
V has unobserved true quality qv , where 0 ≤ qv ≤ 1.                            rating system only through the prior choice, determined here
   Binary ratings that accumulate over time. Products ac-                       by α̂ and β̂. There are two aspects of this choice: the shape
cumulate ratings over time, as users interact with and rate                     of the prior, and its strength relative to observed ratings in
products. To capture this phenomenon, we assume the mar-                        determining the posterior. In our analysis, we fix the prior
ket operates over discrete time periods t, and products can                     shape by relying on market-level data, in a manner inspired
receive ratings at each time period. We assume these ratings                    by empirical Bayesian statistical methods (Robbins 1964).
are binary, i.e., 0 or 1.                                                       Our analysis shows that the strength of the prior then medi-
   Crucially, the true quality of a product should impact the                   ates the tradeoff between efficiency and producer fairness.
rating received; we suppose that a rating for product v is                         Formally, to reason about a prior’s strength, we fix a prior
Bernoulli(qv ), independent of all other randomness in the
                                                                                shape (α̃, β̃), then set (α̂, β̂) = (η α̃, η β̃) for some prior
system. Since ratings accumulate, we use S(v, T ) to denote
                                                                                strength η ≥ 0, providing us with a single-parameter lever
the set of binary ratings that product v ∈ V has received
                                                                                to adjust the strength of the prior. We elaborate on how we
after T timesteps have elapsed.
                                                                                calibrate the prior shape in Section 3.2. Note in particular
   We assume binary ratings as a simplification for ease of
                                                                                that if η = 0, then the platform simply implements sample
exposition; our results qualitatively extend to other settings,
                                                                                averaging to compute the estimated quality of an item.
as we demonstrate in Appendix E. We note that this choice
captures how many platforms elicit ratings (e.g., in any plat-                  2.2   Fixed and Responsive Models
form that uses a thumbs-up/thumbs-down system); in other
platforms with more than two options, ratings often follow a                    We study two models for how customers may interact with
“J-curve” where users mostly only use the extreme options                       items, and items accumulate ratings over time. We begin
(Hu, Zhang, and Pavlou 2009). In Appendix G, we demon-                          with a fixed model in which the products receive ratings at
strate that our results hold even when we assume users’ dis-                    the same rate, regardless of their past ratings. We then study
tribution over ratings is centered or even right-skewed.                        a (more realistic) responsive model in which products with
   Prior-Weighted Rating System. In a prior-weighted rat-                       higher estimated ratings receive ratings more often.
ing system, the prior distribution is a design choice of the                       Fixed model. In the fixed setting, all products stay in the
platform. The system operates by estimating a posterior dis-                    market for the entire time horizon, and each product v ∈ V
tribution of the estimated true quality for each product, given                 accumulates one rating at each timestep t. In this model, the
the prior and the received ratings on that product.                             platform eventually accurately learns the true quality of all
   We primarily consider rating systems where the prior rep-                    products, but has noisy estimates for each finite time t.
resents a Beta distribution.1 Then, in this setting, the plat-                     Responsive model. The responsive model differs from
form chooses the prior Beta distribution parameters (α̂, β̂).2                  the fixed model in two ways: (1) items enter and exit the
                                                                                marketplace over time; (2) items are more likely to receive
We will refer to α̂ + β̂ as the prior strength, as it indicates                 ratings if they have a higher estimated quality in the rating
how strongly the prior is concentrated around its mean (and                     system. These changes in the responsive model capture se-
how many ratings will be needed to “overwhelm”) the prior.                      lection effects (Acemoglu et al. 2022), in which consumers
   After a product receives a rating, the platform updates its                  may react to previous ratings, more often choosing to inter-
estimation posterior distribution for its true quality. In our                  act with (and thus rate) items with more positive ratings in
    1
      This choice is natural, given that ratings are assumed to be
                                                                                the past; these effects are not present in the fixed model.
Bernoulli. The Beta distribution is the conjugate prior for Bernoulli              At each timestep t, only a subset M (t) ⊆ V of the prod-
data, i.e., with this choice of prior, the posterior distribution re-           ucts is available for purchase. A single consumer chooses
mains a Beta distribution but with different parameters.                        a single product v(t) to purchase and provides a rating on
    2
      In a Beta(α, β) distribution, the mean is α+βα
                                                     ; the higher that          that product, and at the end of the timestep, each product
α + β are, the more concentrated the distribution around the mean.              v ∈ M (t) exits the marketplace with some uniform, ex-
A common interpretation of the parameters is that α+β is the num-               ogenous probability and is replaced by another product not
ber of trials (coin flips), and α represents the number of successes.           currently in the market. The customer choice will depend

[Next verbatim source excerpt]

   compromising anonymity...                                                       Bias2 = (E[q̂ηα̃,ηβ̃ (v, t)|qv ] − qv )2                        (6)
                                                                                                "                     #     !2
 (a) If your work uses existing assets, did you cite the cre-                                     η α̃ + R(v, t)
      ators? Yes; we cite the datasets used (Gao et al. 2022;                            = E                        qv − qv                        (7)
                                                                                                  η α̃ + η β̃ + t
      Wan et al. 2019; Wan and McAuley 2018; Hou et al.
                                                                                                                      !2
      2024).                                                                                 η α̃ − η α̃qv − η β̃qv
 (b) Did you mention the license of the assets? Yes, when                                =                                .                        (8)
      we first perform in-depth analysis for each dataset.                                        η α̃ + η β̃ + t
 (c) Did you include any new assets in the supplemental                     Variance can also be broken down accordingly:
      material or as a URL? No, because we did not develop
      any new datasets or assets outside of our simulation
      experiments.                                                          Variance = Var[q̂ηα̃,ηβ̃ (v, t)|qv ]                    (9)
 (d) Did you discuss whether and how consent was ob-                                     "
                                                                                           [U+0010 control character]                                   [U+0011 control character]2
                                                                                                                                   #
      tained from people whose data you’re using/curating?                           = E q̂ηα̃,ηβ̃ (v, t) − E[q̂ηα̃,ηβ̃ (v, t)] qv
      No, because the license for the dataset permits sharing
      and adaptation of the content so long as it is properly                                                                                 (10)
      attributed.                                                                                                           2
                                                                                        E[(R(v, t) − E[R(v, t)]) |qv ]
 (e) Did you discuss whether the data you are using/cu-                               =                                                       (11)
      rating contains personally identifiable information or                                    (η α̃ + η β̃ + t)2
      offensive content? No, because the original authors of                            Var[R(v, t)|qv ]
                                                                                      =                                                       (12)
      the dataset fully anonymized it.                                                  (η α̃ + η β̃ + t)2
  (f) If you are curating or releasing new datasets, did you                              tqv (1 − qv )
      discuss how you intend to make your datasets FAIR?                              =                    .                                  (13)
      N/A                                                                               (η α̃ + η β̃ + t)2
 (g) If you are curating or releasing new datasets, did you                 The expected MSE can be written as follows:
      create a Datasheet for the Dataset? N/A
6. Additionally, if you used crowdsourcing or conducted                                                            !2
   research with human subjects, without compromising                                   η α̃ − η α̃qv − η β̃qv                tqv (1 − qv )
                                                                             MSE =                                      +                      .
   anonymity...                                                                              η α̃ + η β̃ + t                (η α̃ + η β̃ + t)2
\end{ReviewVerbatim}



\paragraph{Lean-elaborated paper signature}
\begin{ReviewVerbatim}
type:
{alpha beta eta q : ℝ} → {t : ℕ} → 0 < alpha → 0 < beta → 0 ≤ eta → 0 < t → 0 ≤ q → q ≤ 1 → Prop

value:
fun {alpha beta eta q : ℝ} {t : ℕ} (halpha : 0 < alpha) (hbeta : 0 < beta) (heta : 0 ≤ eta) (ht : 0 < t) (hq0 : 0 ≤ q)
    (hq1 : q ≤ 1) =>
  MBJG25ProducerFairness.paper_expected_posterior_mean alpha beta eta t q =
      (eta * alpha + ↑t * q) / (eta * alpha + eta * beta + ↑t) ∧
    MBJG25ProducerFairness.paper_conditional_bias alpha beta eta t q =
        MBJG25ProducerFairness.paper_expected_posterior_mean alpha beta eta t q - q ∧
      MBJG25ProducerFairness.paper_conditional_variance alpha beta eta t q =
          ↑t * q * (1 - q) / (eta * alpha + eta * beta + ↑t) ^ 2 ∧
        MBJG25ProducerFairness.paper_conditional_squared_bias alpha beta eta t q =
          MBJG25ProducerFairness.paper_conditional_bias alpha beta eta t q ^ 2
\end{ReviewVerbatim}


\paragraph{Recorded source-to-declaration screening}
\textbf{Verdict:} \texttt{matches}
\\
{\raggedright
\textbf{Reason:} The source specifies binary Bernoulli ratings in discrete timesteps, one rating per product per fixed-\allowbreak{}model timestep, Beta parameters (eta*alpha, eta*beta) with eta >= 0, and its Appendix-\allowbreak{}A bias and variance formulas.\allowbreak{} The target faithfully uses a positive natural review count cast to ℝ and states the corresponding conditional expectation, bias, variance, and squared-\allowbreak{}bias identities.\allowbreak{}
\par}
\reviewmatch{Matches source input}
\noindent\textbf{Reviewer annotation}\par
\reviewerbox

\clearpage
\hypertarget{source-claim-98cc132a8c76d666}{}
\hypertarget{source-section-f10b5f68e4acf3dc}{}
\section*{Source claims}
\subsection*{Review row 1}
\textbf{Source locator:} {\footnotesize\raggedright cited publication, lines 313-\allowbreak{}334\par}
\paragraph{Verbatim source input}
\begin{ReviewVerbatim}
Theorem 3.1. Suppose we have a prior-weighted rating sys-                      iments, is attached in the submission.
tem in the fixed setting with prior parameters (η α̃, η β̃). Fix                  Model calibration. To pick a prior shape (α̃, β̃), we make
product v with true quality qv and consider quality estima-                    use of the empirical Bayes (EB) method (Robbins 1964),
tion after t timesteps.                                                        which provides an approach for computing an estimate of
   Then, as η             ≥ 0 increases, (1) squared bias                      prior parameters across a population. We perform a 60-40
[U+0010 control character]                              [U+0011 control character]2                                              train-test split on the products in the our dataset. We use all
  E[q̂ηα̃,ηβ̃ (v, t)|qv ] − qv     is nondecreasing; (2) vari-                 available ratings data for products in the training set to cali-
ance Var[q̂ηα̃,ηβ̃ (v, t)|qv ] is strictly decreasing in η.                    brate (α̃, β̃) using EB. The remaining 1,331 products (in the
                                                                               test set) are the universe of products V on which the simu-
   Theorem 3.1 establishes the trade-off between fairness                      lation is run. We are interested in investigating how ratings
and efficiency. As the prior strength of a prior-weighted rat-                 in the marketplace change over time, so our marketplace is
ing system increases, it is less likely for two identical prod-                simulated over time; we fix a finite time horizon T , and at
ucts to be rated differently, reducing variance while increas-                 each timestep t, either draw one review for all v ∈ V in the
ing estimation bias.                                                           fixed model, or draw one review for the specific v ∈ V cho-
   The above theorem characterizes error for a given quality                   sen by the consumer at timestep t in the responsive model.
level qv as prior strength changes. We are also interested in                     Fixed model parameters. We set a time horizon of T =
how the true quality of the product influences the efficiency                  50, and run our simulations for four different values of
and fairness metric, i.e., how the platform estimates quality                  η ∈ {0, 1, 10, 1000}. The first two η values correspond to
for different true quality levels. An ideal statistical estimator              prior-weighted rating systems under the sample mean esti-
for quality should have low MSE across all quality levels.                     mator and the baseline empirical Bayes (EB) estimator re-
Our next theorem characterizes how efficiency and unfair-                      spectively. η = 10 corresponds to a setting where the prior
ness vary with true product quality:                                           is difficult to overcome in early timesteps, but eventually is
\end{ReviewVerbatim}


\paragraph{Formalized review target}
The archival source above is not asserted equivalent to this formalized target.
\begin{ReviewVerbatim}
Corrected Theorem 3.1: in the fixed Beta-Bernoulli model, after a positive discrete review count and for nonnegative prior strengths with etaLow <= etaHigh, squared conditional bias is nondecreasing; conditional variance weakly decreases for every true quality q in [0,1], and strictly decreases when 0 < q < 1 and etaLow < etaHigh.
\end{ReviewVerbatim}

\textbf{Archival source anchor:} {\footnotesize\raggedright cited publication, lines 313-\allowbreak{}320\par}
\paragraph{Expanded PaperInterface specification}
\begin{ReviewVerbatim}
∀ {alpha beta q etaLow etaHigh : ℝ} {t : ℕ},
  0 < alpha →
    0 < beta →
      0 < t →
        0 ≤ q →
          q ≤ 1 →
            0 ≤ etaLow →
              etaLow ≤ etaHigh →
                MBJG25ProducerFairness.paper_conditional_squared_bias alpha beta etaLow t q ≤
                    MBJG25ProducerFairness.paper_conditional_squared_bias alpha beta etaHigh t q ∧
                  MBJG25ProducerFairness.paper_conditional_variance alpha beta etaHigh t q ≤
                      MBJG25ProducerFairness.paper_conditional_variance alpha beta etaLow t q ∧
                    (0 < q →
                      q < 1 →
                        etaLow < etaHigh →
                          MBJG25ProducerFairness.paper_conditional_variance alpha beta etaHigh t q <
                            MBJG25ProducerFairness.paper_conditional_variance alpha beta etaLow t q)
\end{ReviewVerbatim}



\paragraph{Lean proof endpoint (built separately)}\begin{ReviewVerbatim}
MBJG25ProducerFairness.paper_theorem3_1_proof
\end{ReviewVerbatim}

\paragraph{Recorded source-to-Spec screening}
\textbf{Verdict:} \texttt{matches}
\\
{\raggedright
\textbf{Reason:} The target expresses weakly increasing squared bias and weakly decreasing variance as prior strength rises, with strict variance decrease restricted to 0<q<1 and etaLow<etaHigh, exactly matching the approved endpoint correction;\allowbreak{} positive discrete t and q∈[0,1] are stated approved model context.\allowbreak{}
\par}
\reviewmatch{Matches formalized target}
\noindent\textbf{Reviewer annotation}\par
\reviewerbox
\clearpage
\hypertarget{source-claim-ad27f52afb94912c}{}
\section*{Review row 2}
\textbf{Source locator:} {\footnotesize\raggedright cited publication, lines 335-\allowbreak{}344\par}
\paragraph{Verbatim source input}
\begin{ReviewVerbatim}
Theorem 3.2. Suppose we have a prior-weighted rating sys-                      shifted significantly by ratings data. η = 1000 corresponds
tem in the fixed setting with prior parameters (η α̃, η β̃). Con-              to a setting where the posterior is essentially equal to the
sider quality estimation after t timesteps.                                    prior throughout the entire simulation, i.e., ratings data is
   Then, as a function of true product quality, squared bias                   ignored. We ran all of our experiments (for both the fixed
[U+0010 control character]                              [U+0011 control character]2                                              and responsive models) on an internally-provided comput-
  E[q̂ηα̃,ηβ̃ (v, t)|qv ] − qv    is convex with a global mini-                ing cluster, where the compute power required for our simu-
mum at α̃+α̃
                , while variance Var[q̂ηα̃,ηβ̃ (v, t)|qv ] is concave          lations was negligible.
             β̃
with a global maximum at 1/2.
\end{ReviewVerbatim}



\paragraph{Expanded PaperInterface specification}
\begin{ReviewVerbatim}
∀ {alpha beta eta : ℝ} {t : ℕ},
  0 < alpha →
    0 < beta →
      0 ≤ eta →
        0 < t →
          (∀ (x y lam : ℝ),
              0 ≤ x →
                x ≤ 1 →
                  0 ≤ y →
                    y ≤ 1 →
                      0 ≤ lam →
                        lam ≤ 1 →
                          MBJG25ProducerFairness.paper_conditional_squared_bias alpha beta eta t
                              (lam * x + (1 - lam) * y) ≤
                            lam * MBJG25ProducerFairness.paper_conditional_squared_bias alpha beta eta t x +
                              (1 - lam) * MBJG25ProducerFairness.paper_conditional_squared_bias alpha beta eta t y) ∧
            (0 ≤ alpha / (alpha + beta) ∧
                alpha / (alpha + beta) ≤ 1 ∧
                  ∀ (q : ℝ),
                    0 ≤ q →
                      q ≤ 1 →
                        MBJG25ProducerFairness.paper_conditional_squared_bias alpha beta eta t
                            (alpha / (alpha + beta)) ≤
                          MBJG25ProducerFairness.paper_conditional_squared_bias alpha beta eta t q) ∧
              (∀ (x y lam : ℝ),
                  0 ≤ x →
                    x ≤ 1 →
                      0 ≤ y →
                        y ≤ 1 →
                          0 ≤ lam →
                            lam ≤ 1 →
                              lam * MBJG25ProducerFairness.paper_conditional_variance alpha beta eta t x +
                                  (1 - lam) * MBJG25ProducerFairness.paper_conditional_variance alpha beta eta t y ≤
                                MBJG25ProducerFairness.paper_conditional_variance alpha beta eta t
                                  (lam * x + (1 - lam) * y)) ∧
                ∀ (q : ℝ),
                  0 ≤ q →
                    q ≤ 1 →
                      MBJG25ProducerFairness.paper_conditional_variance alpha beta eta t q ≤
                        MBJG25ProducerFairness.paper_conditional_variance alpha beta eta t (1 / 2)
\end{ReviewVerbatim}



\paragraph{Lean proof endpoint (built separately)}\begin{ReviewVerbatim}
MBJG25ProducerFairness.paper_theorem3_2_proof
\end{ReviewVerbatim}

\paragraph{Recorded source-to-Spec screening}
\textbf{Verdict:} \texttt{matches}
\\
{\raggedright
\textbf{Reason:} The target formalizes the source claims on true-\allowbreak{}quality domain [0,1]:\allowbreak{} squared conditional bias is convex with global minimum at alpha/\allowbreak{}(alpha+beta), and conditional variance is concave with global maximum at 1/\allowbreak{}2;\allowbreak{} positive discrete t is covered by the approved model context.\allowbreak{}
\par}
\reviewmatch{Matches source input}
\noindent\textbf{Reviewer annotation}\par
\reviewerbox

\end{document}
